\section{Por qué el Generative Modeling puede convertirse en Unsupervised Learning}

El aprendizaje supervisado depende de etiquetas $(x,y)$, mientras que gran parte de los datos de internet solo tiene $x$. El modelado generativo construye la señal de entrenamiento al predecir los propios datos, por lo que puede aprender representation y behavior sin etiquetas de tarea creadas por personas. La cuestión central no es que haya «cero etiquetas», sino si el objective autosupervisado obliga al modelo a capturar estructura útil para las tareas posteriores.

\subsection{El cuello de botella de las etiquetas y la expansión de tareas}

En la clase, \term{unsupervised learning} (aprendizaje no supervisado) designa en general el aprendizaje de representaciones que no depende de etiquetas de tarea creadas por personas; en el contexto actual suele llamarse con más precisión \term{self-supervised learning}, porque el target se construye automáticamente a partir de los datos originales. La next-token label de un modelo de lenguaje proviene del propio texto.

\lecturefigure{slide-10-unsupervised-learning.jpg}{La demanda de datos de las tareas supervisadas lleva a los modelos a buscar señales de entrenamiento que puedan generarse automáticamente a partir de datos sin procesar.}{Video oficial 09:09.}

\begin{knowledgebox}{Lectura de la figura: la diferencia entre ambas rutas está en cómo se producen las etiquetas}
El supervised learning requiere task-specific labels, que son costosas y de cobertura limitada; la ruta unsupervised/self-supervised aprende estructura general a partir de raw data a gran escala y luego la transfiere mediante prompting o fine-tuning. Esto no elimina las human choices: la recolección y el filtrado de datos, el tokenizer, el objective y la evaluation siguen siendo decisiones de diseño humanas.
\end{knowledgebox}

\lecturefigure{slide-11-unlabeled-data.jpg}{Internet ofrece cantidades enormes de datos sin etiquetar, pero la escala no equivale a calidad ni a autorización.}{Video oficial 10:38.}

\begin{knowledgebox}{Lectura de la figura: el ícono del planeta representa la escala y también la gobernanza}
Los web data a gran escala abarcan diferencias de idioma, código, imágenes y cultura, y aportan una signal amplia para el preentrenamiento; al mismo tiempo contienen duplication, bias, private data, copyright y contamination. La clase se centra en «si existe un gran acervo de datos»; en la práctica de ingeniería también hay que responder por el origen, la licencia, la limpieza, la deduplicación, la retención y la eliminación.
\end{knowledgebox}

\subsection{Analysis by synthesis y el autoregressive objective}

La intuición de \term{analysis by synthesis} es la siguiente: si un modelo puede generar muestras que coinciden con la distribución de los datos, debe haber aprendido parte de su estructura latente. El autoregressive objective maximiza
\[
\mathcal{L}(\theta)=\sum_{t=1}^{T}\log p_\theta(x_t\mid x_{<t}).
\]
Cada posición produce una señal de entrenamiento, lo que lo hace adecuado para escalar (scale); pero el objective solo exige predecir tokens y no garantiza que la representation interna corresponda uno a uno con los conceptos humanos.

\lecturefigure{slide-12-autoregressive-analysis-by-synthesis.jpg}{La autoregressive generation aprende estructura al predecir los datos; es una forma de implementar analysis by synthesis.}{Video oficial 11:11.}

\begin{knowledgebox}{Lectura de la figura: la generación es la interfaz de entrenamiento; la representation es un subproducto}
Para predecir el siguiente token, el modelo aprovecha la sintaxis, el tema, las entidades, el sentimiento y las regularidades del mundo; estos factores pueden formar direcciones legibles en los hidden states. Pero el modelo también puede aprovechar un surface shortcut. Para juzgar si una representation es útil se necesitan probing, intervention y transfer, no solo una loss baja.
\end{knowledgebox}

\subsection{Sentiment neuron: representaciones emergentes y los límites de la interpretación}

Radford et al. encontraron, en un character-level language model, una unidad fuertemente correlacionada con el sentimiento de las reseñas (review sentiment). Este resultado muestra que un objective generativo puede aprender features útiles para clasificación, y también recuerda que «una neuron» suele ser un fenómeno propio de un modelo y un método de análisis específicos; no debe extrapolarse a la idea de que todos los conceptos se codifican en una sola neurona.

\lecturefigure{slide-13-sentiment-neuron.jpg}{En un modelo generativo no supervisado aparece una dirección de representación relacionada con el sentimiento.}{Video oficial 13:00; Radford et al. 2017.}

\begin{knowledgebox}{Lectura de la figura: a la izquierda, la distribución; a la derecha, la trayectoria de activación}
El lado izquierdo compara la diferencia de distribución de las reviews positivas y negativas en la neuron activation; el lado derecho muestra la posición en el texto y la variación de la activation. La correlación respalda que esta unidad transporta una sentiment signal, pero por sí sola no demuestra que sea el único mecanismo causal. La causalidad y la generalización pueden ponerse a prueba con clamping, ablation o classifier transfer.
\end{knowledgebox}

\teachervoice{El ponente usa la sentiment neuron para responder «por qué generar texto ayuda a las tareas posteriores»: para hacer synthesis, el modelo posiblemente primero haga analysis. La advertencia de la clase es que se trata de una posibilidad, no de una garantía; la representation quality todavía debe validarse con evidencia en las tareas posteriores.}

\subsection{Resumen de la sección}

El preentrenamiento generativo usa los propios datos para construir supervisión densa y puede aprender representaciones transferibles a partir de la escala de datos sin etiquetar. Su ventaja proviene de la cobertura de la señal; su riesgo, de la distancia que sigue existiendo entre el objective y las tareas reales.
